\section{The one-point specialization}
\label{sec:BA-one-point-specialization}

The general construction becomes a vector--matrix--vector system
when $N=1$. We give its elementary-mode notation, specialize the
higher-flow phases and explain its constrained-KP interpretation.
The curve remains arbitrary throughout this section.

\subsection{Elementary modes and their canonical equations}
\label{sec:BA-one-point}

For $N=1$ put $R_1=P_\infty$, $s_1=s$ and $e_1=1$.
The divisors have degree $g$ and there is a single ordinary/dual pair
$\psi,\psi^\vee$. Choose reference normalization
$\psi(P,0)=\psi^\vee(P,0)=1$, and subtract the constant term of
each $h$ at $P_\infty$. Thus $a_{L,1}=a_{M,1}=0$.
This member has two elementary vectors and a pair-mode matrix.

The combined index and number of canonical pairs are
\begin{equation}\label{eq:BA-mode-counts}
 K=L+M,\qquad p=K+LM.
\end{equation}
We use single indices for the $K$ elementary pairs and double indices
for the $LM$ cross-group pair composites. Their dictionary with the
$H_0$ modes of \cref{eq:canonical-phase-gauge} is
\begin{equation}\label{eq:BA-mode-dictionary}
 \xi_i=\tilde w_{1i},\qquad
 \xi_{L+m}=\tilde v_{m1},\qquad
 \xi_{i,L+m}=\tilde u_{im},
 \quad1\leq i\leq L,\quad1\leq m\leq M,
\end{equation}
with the same rule for barred partners. Put
\begin{equation}\label{eq:BA-phase-dictionary}
 b_i=2b_i^L+b_1^N,\qquad b_{L+m}=2b_m^M+b_1^N,
 \quad1\leq i\leq L,\quad1\leq m\leq M.
\end{equation}
The inherited brackets satisfy $\{\xi_k,\bar\xi_r\}=-i\delta_{kr}$, with the corresponding rule for pair modes, and
\begin{equation}\label{eq:BA-H0-flat}
 H_0=\sum_{k=1}^{K}b_k\xi_k\bar{\xi}_k
 +\sqrt2\sum_{i=1}^{L}\sum_{j=L+1}^{K}
 (\xi_i\xi_j\xi_{ij}+\bar{\xi}_i\bar{\xi}_j\bar{\xi}_{ij}).
\end{equation}


The equivalent one-point theta and sigma formulas, with
$\Xi=s^{-1}+O(s)$ and
$\Xi_\sigma=\Xi-\boldsymbol U^{\trans}\varkappa A(P)$, are
\begin{equation}\label{eq:general-BA-theta}
\begin{aligned}
 \psi(P,t)
 &=e^{-it\Xi(P)}
 \tfrac{\theta(z(t)+A(P))\theta(z_0)}
 {\theta(z(t))\theta(z_0+A(P))}
 =e^{-it\Xi_\sigma(P)}
 \tfrac{\sigma(z(t)+A(P))\sigma(z_0)}
 {\sigma(z(t))\sigma(z_0+A(P))},\\[4pt]
 \psi^\vee(P,t)
 &=e^{it\Xi(P)}
 \tfrac{\theta(z(t)-A(P))\theta(z_0)}
 {\theta(z(t))\theta(z_0-A(P))}
 =e^{it\Xi_\sigma(P)}
 \tfrac{\sigma(z(t)-A(P))\sigma(z_0)}
 {\sigma(z(t))\sigma(z_0-A(P))}.
\end{aligned}
\end{equation}
In this notation the unphased BA modes are
\begin{equation}\label{eq:general-elementary-modes}
\begin{aligned}
 \xi_i&=c_i \psi^\vee(Q_i),&\bar{\xi}_i&=\kappa_i \psi(Q_i)/c_i,
 &&1\leq i\leq L,\\
 \xi_j&=c_j \psi(Q_j),&\bar{\xi}_j&=\kappa_j \psi^\vee(Q_j)/c_j,
 &&L+1\leq j\leq K,
\end{aligned}
\end{equation}
and
\begin{equation}\label{eq:general-pair-modes}
 \xi_{ij}=\frac{\sqrt2\kappa_i\kappa_j}{c_ic_j}\Omega(Q_i,Q_j),
 \qquad \bar{\xi}_{ij}=-\sqrt2c_ic_j\Omega(Q_j,Q_i),
 \quad1\leq i\leq L<j\leq K.
\end{equation}

An \emph{elementary mode} has the shift of one curve-point lift
relative to $P_\infty$. A \emph{pair composite} has the difference
of the two lifts; in the hyperelliptic sheet convention this becomes
the sum of the elementary shifts. The terms specify the shift data,
not a factorization at a common $z$. For $N>1$, the basis sections
instead have shifts $A(P)-A(R_k)$; all sections share the same curve
and kernel, but generally no single base point makes both arrays
one-point elementary vectors.

The raw phase values and equations reduce to
\begin{equation}\label{eq:one-point-phase-values}
 b_i=-h_M(Q_i),\quad1\leq i\leq L,\qquad
 b_j=h_L(Q_j),\quad L+1\leq j\leq K,
\end{equation}
\begin{equation}\label{eq:general-network-rates}
\begin{aligned}
 i\dot\xi_i&=b_i\xi_i+\sqrt2\sum_{j=L+1}^{K}\bar{\xi}_j\bar{\xi}_{ij},&
 i\dot{\bar{\xi}}_i&=-b_i\bar{\xi}_i-\sqrt2\sum_{j=L+1}^{K}\xi_j\xi_{ij},
 &&1\leq i\leq L,\\
 i\dot\xi_j&=b_j\xi_j+\sqrt2\sum_{i=1}^{L}\bar{\xi}_i\bar{\xi}_{ij},&
 i\dot{\bar{\xi}}_j&=-b_j\bar{\xi}_j-\sqrt2\sum_{i=1}^{L}\xi_i\xi_{ij},
 &&L+1\leq j\leq K,\\
 i\dot\xi_{ij}&=\sqrt2\bar{\xi}_i\bar{\xi}_j,&
 i\dot{\bar{\xi}}_{ij}&=-\sqrt2\xi_i\xi_j,
 &&1\leq i\leq L<j\leq K.
\end{aligned}
\end{equation}
These are Hamilton's equations for \cref{eq:BA-H0-flat}.
For this phase gauge, write $\mathsf A_L=\diag(b_1,\ldots,b_L)$,
$\mathsf A_M=\diag(b_{L+1},\ldots,b_K)$.
An optional scalar spectral origin $\rho$ gives
\begin{equation}\label{eq:BA-H0-pencil}
\begin{aligned}
 \mathsf L_0(\nu)&=\begin{pmatrix}
 \rho\Id_L-\mathsf A_L+\nu\Id_L&\sqrt2\boldsymbol\xi_{LM}&\bar{\boldsymbol\xi}_L\\
 \sqrt2\bar{\boldsymbol\xi}_{LM}^{\trans}&\rho\Id_M-\mathsf A_M-\nu\Id_M&\boldsymbol\xi_M\\
 -\boldsymbol\xi_L^{\trans}&-\bar{\boldsymbol\xi}_M^{\trans}&\rho
 \end{pmatrix},\\
 \mathsf T_0(\nu)&=\begin{pmatrix}
 0&0&\bar{\boldsymbol\xi}_L\\0&0&-\boldsymbol\xi_M\\
 -\boldsymbol\xi_L^{\trans}&\bar{\boldsymbol\xi}_M^{\trans}&-\nu
 \end{pmatrix}.
\end{aligned}
\end{equation}
Here $\boldsymbol\xi_L,\boldsymbol\xi_M$ are the column vectors of
the two elementary groups and $\boldsymbol\xi_{LM}$ is their pair
matrix. This is the one-point form of the canonical pencil after
its phase gauge; $\rho$ only shifts the spectral ordinate.

\subsection{Higher-flow phases in the one-point notation}
\label{sec:BA-one-point-higher-flows}

We specialize \cref{thm:BA-higher-flows} using the modes and pencil
just defined. This gives a modal-power expression for the finite
phase constants of a higher Hamiltonian.

For $N=1$, choose the ordinary gap complement at $P_\infty$.
Then $p_1=\sum_{\alpha=1}^{g}v_\alpha s^{-n_\alpha}$ for its
Weierstrass gaps $n_\alpha$, and $\varphi_1^N=0$.
The elementary phase rates are
$b_i^{(r,\nu_0)}=-\varphi_i^L$,
$b_{L+j}^{(r,\nu_0)}=\varphi_j^M$; pair phases are the opposite
sums of these rates. For example, in the one-point $H_0$ pencil
\cref{eq:BA-H0-pencil}, with optional spectral origin $\rho$, the
fourth trace has $q=\lambda^3/\nu$. Put $\epsilon_a=\rho-b_a$,
$\chi_a=1$ in the first ordinary group and $\chi_a=-1$ in the second.
Write $I_a=\xi_a\bar\xi_a$, $1\leq a\leq K$, and
$I_{ij}=\xi_{ij}\bar\xi_{ij}$, $1\leq i\leq L<j\leq K$.
For the conserved power combinations
$C_i=I_i-\sum_{j=L+1}^{K}I_{ij}$ and
$C_j=I_j-\sum_{i=1}^{L}I_{ij}$, perturbation at infinity gives
\[
 \lambda=\chi_a\nu+\epsilon_a-\chi_aC_a/\nu+O(\nu^{-2}),\qquad
 \operatorname{reg}_{Q_a}q=3\chi_a(\epsilon_a^2-C_a).
\]
The phase is this constant plus the finite part of $\mathfrak f$,
with the elementary sign just specified. Thus the one-point
modal-power interpretation is retained inside the multipoint theorem.
